% ---------------------------------------------------------------------
%  CHƯƠNG 6 — KẾT LUẬN VÀ HƯỚNG PHÁT TRIỂN
% ---------------------------------------------------------------------
\chapter{Kết luận và hướng phát triển}
\label{ch:ketluan}

\begin{muctieu}
Chương này tổng kết mức độ hoàn thành của các mục tiêu nghiên cứu, đúc kết các
bài học kinh nghiệm quý giá trong quá trình phát triển hệ thống AI y tế và đề
xuất các hướng phát triển mở rộng tiếp theo của đề tài.
\end{muctieu}

\section{Kết luận}

Đồ án đã xây dựng thành công \textbf{Hệ thống trợ lý sàng lọc sức khỏe và phân
tầng nguy cơ y tế} tích hợp trí tuệ nhân tạo lai, đi trọn vòng đời từ phân tích
yêu cầu, thiết kế cơ sở tri thức, cài đặt các dịch vụ suy diễn đến triển khai
thực nghiệm đo lường định lượng. Hệ thống đã hoàn thành cả 4 mục tiêu nghiên
cứu đã xác lập:

\begin{itemize}[leftmargin=1.6em, itemsep=2pt]
  \item \textbf{Về mặt khoa học}: Đã hiện thực hóa kiến trúc Hybrid AI bền vững,
        trong đó động cơ quy tắc tất định niêm phong mọi quyết định an toàn
        (\textbf{88,5\%} xử trí an toàn lâm sàng trên 26 ca thực nghiệm), còn
        mô hình ngôn ngữ lớn đảm nhận lớp tương tác và diễn giải (\textbf{100\%}
        đạt chuẩn kiểm duyệt, độ phủ \textbf{5/5} danh mục hướng dẫn trên mỗi ca).
  \item \textbf{Về mặt kỹ thuật}: Đã triển khai ứng dụng web Laravel hoàn chỉnh
        với kiến trúc hàng đợi bất đồng bộ (\code{AiJob}), cơ chế cứu hộ đồng bộ
        khi worker bị treo và hệ thống hàng đợi thẩm duyệt của dược sĩ, được kiểm
        chứng bằng bộ kiểm thử tự động 298 ca (965 phép khẳng định) đều vượt qua.
  \item \textbf{Về mặt xã hội}: Giải pháp mang lại công cụ định hướng y tế ban
        đầu tiện lợi, kịp thời và ân cần cho người dân, đồng thời góp phần giảm
        tải cho hệ thống y tế cơ sở trong các bệnh cảnh thông thường.
\end{itemize}

\section{Bài học rút ra}

\begin{enumerate}[leftmargin=1.6em, itemsep=2pt]
  \item \textbf{Trong y tế, kiểm soát đầu ra quan trọng hơn sức mạnh mô hình}:
        Điều quyết định độ an toàn của hệ thống không phải là chọn mô hình ngôn
        ngữ lớn mạnh nhất, mà là thiết kế các lớp kiểm duyệt tại tầng mã nguồn
        (\code{sanitize}, \code{severerRule}). Với chỉ vài dòng lệnh PHP kiểm
        tra \emph{«tối thiểu 2 danh mục»} và phép so sánh độ nghiêm trọng, toàn
        bộ nguy cơ ảo giác y khoa nghiêm trọng đã được chặn đứng tận gốc.
  \item \textbf{Cơ sở tri thức là linh hồn của hệ chuyên gia}: Cả ba trường hợp
        phân tầng thiếu hụt trong thực nghiệm đều xuất phát từ khoảng trống cơ
        sở tri thức (thiếu cụm từ FAST của đột quỵ, thiếu dấu hiệu mắt trũng
        trong mất nước, thiếu nhánh xét bệnh nền đa thuốc), chứ không phải do
        thuật toán suy diễn sai. Đối chứng ở Mục~4.5 xác nhận điều này: bật hay
        tắt lớp AI cho kết quả y hệt trên cả 26 ca. Nếu làm lại, nhóm sẽ dành
        thời gian đầu tiên để hoàn thiện từ điển lâm sàng cùng thầy thuốc chuyên
        môn trước khi tối ưu hoá các lớp AI.
  \item \textbf{Kiến trúc bất đồng bộ giải quyết mâu thuẫn thời gian}: Mâu thuẫn
        cố hữu giữa việc AI tạo sinh cần hàng chục giây để suy nghĩ và yêu cầu
        trải nghiệm web phải phản hồi trong vài giây chỉ có thể hóa giải bằng
        cách tách lớp quyết định (đồng bộ, tức thì) và lớp diễn giải (bất đồng
        bộ, thăm dò). Đây là bài học nền tảng cho mọi ứng dụng tích hợp LLM.
\end{enumerate}

\section{Hướng phát triển}

\begin{itemize}[leftmargin=1.6em, itemsep=2pt]
  \item \textbf{Mở rộng tri thức cờ đỏ và bối cảnh bệnh nền}: Bổ sung các từ khóa
        thang điểm FAST (méo miệng, yếu tay chân, nói khó), dấu hiệu mất nước
        tiến triển (mắt trũng, tiểu ít), và một nhánh luật nâng mức khi hồ sơ có
        bệnh nền đang dùng nhiều thuốc — nhằm khắc phục ba trường hợp phân tầng
        thiếu hụt đã phát hiện, nâng tỷ lệ xử trí an toàn tiến về 100\%.
  \item \textbf{Tích hợp nhận diện giọng nói tiếng Việt (STT)}: Cho phép người
        cao tuổi hoặc người bệnh nặng chỉ cần nói trực tiếp vào thiết bị di động
        thay vì gõ chữ, mở rộng khả năng tiếp cận tới nhóm dân cư dễ tổn thương nhất.
  \item \textbf{Nâng cấp mô hình trích xuất sự kiện tiếng Việt chuyên sâu}: Thử
        nghiệm các mô hình ngôn ngữ tiền huấn luyện dành riêng cho tiếng Việt
        như PhoBERT \cite{Nguyen2020} để nâng cao khả năng nhận diện triệu chứng
        với lối diễn đạt dân gian, phương ngữ vùng miền.
  \item \textbf{Liên thông hồ sơ bệnh án điện tử (EHR)}: Kết nối tự động với lịch
        sử khám chữa bệnh và đơn thuốc điện tử của người bệnh tại các cơ sở y tế,
        giúp hệ thống có thêm bối cảnh tiền sử và nâng cao độ chính xác phân tầng.
  \item \textbf{Triển khai quy mô sản xuất (Production Deployment)}: Đóng gói hệ
        thống thành kiến trúc vi dịch vụ (Microservices), áp dụng bộ cân bằng tải
        và giám sát nhật ký y tế tập trung để phục vụ đồng thời hàng ngàn người
        dùng tại nhiều nhà thuốc và phòng khám trên toàn quốc.
\end{itemize}
